\documentclass[11pt,a4paper]{article}
\usepackage[margin=1.8cm]{geometry}
\usepackage{amsmath,amssymb,tabularx,array,needspace,fontspec,enumitem,cancel}
\newcommand{\tb}{\cancel{||||}\,}
\setmainfont{FreeSerif}
\newfontfamily\urdufont[Script=Arabic]{FreeSerif}
\TeXXeTstate=1
\newcommand{\urdu}[1]{{\urdufont\beginR #1\endR}}
\setlength{\parindent}{0pt}
\setlength{\parskip}{0pt}
\renewcommand{\arraystretch}{1.15}
\pagestyle{plain}

% \mcq{number}{question}{A}{B}{C}{D}
\newcommand{\mcq}[6]{%
  \Needspace{4\baselineskip}%
  \par\medskip\noindent\textbf{#1.}\ #2\par\nopagebreak
  \noindent\begin{tabularx}{\linewidth}{@{\hspace{1.5em}}*{4}{>{\raggedright\arraybackslash}X}@{}}
    (A)~#3 & (B)~#4 & (C)~#5 & (D)~#6
  \end{tabularx}\par}

\begin{document}

\begin{center}
  {\Huge\bfseries The Hope Science Academy}\\[3pt]
  {\small Excellence is our standard not goal}\\[2pt]
  {\small 247-E-1, (Opposite LDA School) Johar Town, Lahore. PH: 0300-8194789, 0322-4157001}
\end{center}

\medskip
\noindent
\begin{tabularx}{\linewidth}{|X|X|X|X|}
\hline
\textbf{Student Name:} & \textbf{Roll Num:} & \textbf{Class:} 9th & \textbf{Paper Code:} \\ \hline
\textbf{Subject Name:} Mathematics. & \textbf{Time Allowed:} & \textbf{Total Marks:} 543 & \textbf{Exam Date:} 09-Oct-2026 \\ \hline
\multicolumn{4}{|l|}{\textbf{Exam Syllabus:} FULL BOOK} \\ \hline
\end{tabularx}

\bigskip
\noindent\textbf{Q1. Choose the correct answer.}\hfill\textbf{1$\times$543=543}

\mcq{1}{The Sumerians used which numeral system:}{10Base 10}{60Base 60}{2Base 2}{12Base 12}
\mcq{2}{The concept of zero was first developed by:}{Egyptians}{Sumerians}{Indians}{Greeks}
\mcq{3}{Modern number system like hexadecimal are based on:}{10Base 10}{16Base 16}{2Base 2}{60Base 60}
\mcq{4}{Which of these is an irrational number:}{0.333...}{3.14159}{1.5}{2332}
\mcq{5}{A recurring decimal is always a:}{Rational number}{Irrational number}{Integer}{Whole number}
\mcq{6}{Which number is rational?}{$\sqrt{3}$}{$\sqrt{16}$}{$\sqrt{5}$}{$\sqrt{7}$}
\mcq{7}{The decimal expansion of $\dfrac{1}{7}$ :}{Terminating}{Non-terminating andnon-recurring}{Non-terminating and recurring}{None of these}
\mcq{8}{Which property is illustrated by a+b = b + a:}{Associative}{Identity}{Distributive}{Commutative}
\mcq{9}{The additive property for real number is:}{1}{0}{-1}{None of these}
\mcq{10}{Which property states a(b+c)=ab+ac:}{Associative}{Distributive}{Inverse}{Closure}
\mcq{11}{The multiplicative inverse of 5 is:}{-5}{$\dfrac{1}{5}$}{0}{1}
\mcq{12}{If a>b and b>c, then a>c is called ....... property:}{Reflexive}{Symmetric}{Transitive}{Trichotomy}
\mcq{13}{Which of the following is not an integer:}{-3}{0}{4.5}{7}
\mcq{14}{Which of the following is an example of a rational number?}{$\pi$}{$\dfrac{3}{4}$}{$\sqrt{2}$}{e}
\mcq{15}{What is the example of a recurring decimal:}{3.14159265 ...}{0.3333 ...}{2.718281828 ...}{0.49835555}
\mcq{16}{Which number system does -3 belong to:}{Natural number}{Whole number}{Integers}{Irrational number}
\mcq{17}{Which of the following is not a real number:}{5}{-2}{$\sqrt{7}$}{0}
\mcq{18}{Which number set the number 0 belong to:}{Whole number}{Natural number}{Irrational number}{None of these}
\mcq{19}{Which of the following is an example of an infinite non-recurring decimal?}{0.25}{0.333...}{5.75}{$\pi$}
\mcq{20}{$0.\overline{6}$ is equivalent to:}{$\dfrac{2}{3}$}{$\dfrac{1}{6}$}{$\dfrac{3}{5}$}{$\dfrac{6}{99}$}
\mcq{21}{Which is a binomial surd?}{$\sqrt{3}$}{$\sqrt{3}+\sqrt{5}$}{$2\sqrt{3}$}{$\sqrt{9}$}
\mcq{22}{The conjugate of $\sqrt{7}-\sqrt{2}$ :}{$\sqrt{7}+\sqrt{2}$}{$\sqrt{7}-\sqrt{2}$}{$-\sqrt{7}+\sqrt{2}$}{$-\sqrt{7}-\sqrt{2}$}
\mcq{23}{Rationalize the denominator $\dfrac{1}{\sqrt{5}}$ :}{$\sqrt{5}$}{$\dfrac{\sqrt{5}}{5}$}{$\dfrac{1}{5}$}{$5\sqrt{5}$}
\mcq{24}{Which is not a surd?4.\urdu{درج ذیل میں سے کون سے مقادیر اصم نہیں ہے؟}}{$\sqrt{8}$}{$\sqrt{12}$}{$\sqrt{16}$}{$\sqrt{18}$}
\mcq{25}{The solution of 2x+5=15 is:}{x=3}{x=10}{x=7.5}{x=5}
\mcq{26}{If $\sqrt{x}=5$ :}{10}{25}{5}{15}
\mcq{27}{The value of $\dfrac{\sqrt{12}\times\sqrt{3}}{\sqrt{4}}$ :}{3}{6}{9}{12}
\mcq{28}{If a=3, b=4, then $\sqrt{a^{2}+b^{2}}$ :}{5}{7}{12}{25}
\mcq{29}{Which is equivalent to $\dfrac{\sqrt{3}+\sqrt{2}}{\sqrt{3}-\sqrt{2}}$ = ?:}{$5+2\sqrt{6}$}{1}{$\sqrt{6}$}{$5-2\sqrt{6}$}
\mcq{30}{If $x=1+\sqrt{2}$, then x$^{2}$=?}{$2\sqrt{2}$}{$1+2\sqrt{2}$}{3}{$3+2\sqrt{2}$}
\mcq{31}{If a+b=10 and ab=21, then a$^{2}$+b$^{2}$=:}{58}{100}{42}{121}
\mcq{32}{The solution of 2$^{x}$=16 is:}{6}{8}{2}{4}
\mcq{33}{What is the square root of 16:}{4}{-4}{Both a and b}{None of these}
\mcq{34}{The value of $\sqrt{25}$ :}{5}{-5}{Both a and b}{None of these}
\mcq{35}{Convert 25$^{\circ}$C to Fahrenheit using $F=\dfrac{9}{5}C+32$:}{77$^{\circ}$F}{68$^{\circ}$F}{90$^{\circ}$F}{45$^{\circ}$F}
\mcq{36}{A profit of 20\% on cost price of Rs. 500 means the selling price is:}{Rs 600}{Rs 400}{Rs 550}{Rs 700}
\mcq{37}{Compound interest on Rs 2000 at 10\% for 2 years is:}{Rs 420}{Rs 200}{Rs 400}{Rs 240}
\mcq{38}{If three consecutive integers sum to 51, the largest number is:}{15}{16}{17}{18}
\mcq{39}{A rectangle with sides $\sqrt{3}$ :}{$\sqrt{15}$}{6}{$3\sqrt{3}$}{12}
\mcq{40}{If 5 pens cost Rs 45, the cost of 8 pens is:}{100}{60}{90}{72}
\mcq{41}{A number is increased by 20\% becomes 60. The original number is:}{50}{48}{60}{72}
\mcq{42}{Two numbers are in the ratio 3:4. Their sum is 35. The numbers are:}{15,20}{12,16}{9,12}{18,24}
\mcq{43}{A train travels 300 km in 5 hours, Its speed is:}{50 km/h}{60 km/h}{45 km/h}{65 km/h}
\mcq{44}{The area of square with side $\sqrt{5}$ :}{5}{10}{25}{20}
\mcq{45}{The sum of $\sqrt{0.25}+\sqrt{0.04}$ :}{0.7}{0.5}{0.9}{1.0}
\mcq{46}{The temperature 100$^{\circ}$C is equal to how many Kelvin:}{273 K}{373 K}{473 K}{173 K}
\mcq{47}{$\sqrt{7}$ :}{Integer}{Rational number}{Irrational number}{Natural number}
\mcq{48}{$\pi$ and e are:}{Natural number}{Integers}{Rational numbers}{Irrational numbers}
\mcq{49}{If n is not a perfect square, then $\sqrt{n}$ :}{\urdu{ناطق عدد} Rational number}{\urdu{قدرتی عدد} Natural number}{\urdu{صحیح عدد} Integer}{\urdu{غیر ناطق عدد} Irrational number}

\mcq{50}{$\sqrt{3}+\sqrt{5}$ :}{Whole number}{Integer}{Rational number}{Irrational number}
\mcq{51}{For all $x\in R$, $x=x$ :}{\urdu{عکسی خاصیت} Reflexive property}{\urdu{متعدیت} Transitive number}{\urdu{تشاکل} Symmetric property}{\urdu{ثلاثی خاصیت} Trichotomy property}
\mcq{52}{Let $a,b,c\in R$, then $a>b$ and $b>c\Rightarrow a>c$ is called ...... property:}{Trichotomy}{Transitive}{Additive}{Multiplicative}
\mcq{53}{$2^{x}\times 8^{x}=64$ then x=}{$\dfrac{3}{2}$}{$\dfrac{3}{4}$}{$\dfrac{5}{6}$}{$\dfrac{2}{3}$}
\mcq{54}{Let $a,b\in R$, then $a=b$ and $b=a$ is called ...... property:}{Reflexive}{Symmetric}{Transitive}{Additive}
\mcq{55}{$\sqrt{75}+\sqrt{27}$ = :}{$\sqrt{102}$}{$9\sqrt{3}$}{$5\sqrt{3}$}{$8\sqrt{3}$}
\mcq{56}{The product of $(3+\sqrt{5})(3-\sqrt{5})$ :}{\urdu{مفرد عدد} Prime number}{\urdu{طاق عدد} Odd number}{\urdu{غیر ناطق عدد} Irrational number}{\urdu{ناطق عدد} Rational number}
\mcq{57}{Convert 45,000 to scientific notation:}{4.5 x 10$^{3}$}{4.5 x 10$^{4}$}{45 x 10$^{3}$}{0.45 x 10$^{5}$}
\mcq{58}{The scientific notation of 0.000078 is:}{$7.8\times 10^{-5}$}{$78\times 10^{-6}$}{$0.78\times 10^{-4}$}{$7.8\times 10^{5}$}
\mcq{59}{Which is not in scientific notation:}{$3.4\times 10^{6}$}{$0.5\times 10^{3}$}{$6.02\times 10^{23}$}{$1.2\times 10^{-4}$}
\mcq{60}{Convert $9.01\times 10^{5}$ :}{9010}{90100}{901000}{901}
\mcq{61}{The coefficient in $6.022\times 10^{23}$ is:}{6}{6.022}{$10^{23}$}{23}
\mcq{62}{$73\times 10^{3}$ in scientific notation is:}{$7.3\times 10^{2}$}{$7.3\times 10^{3}$}{$7.3\times 10^{4}$}{$73\times 10^{3}$}
\mcq{63}{Convert 6.7779 x 10$^{3}$km to standard form:}{677.76 km}{67.779 km}{67779 km}{6777.9 km}
\mcq{64}{The exponential form of $\log_{5}25=2$ :}{25$^{2}$=5}{2$^{5}$=25}{5$^{2}$=25}{25$^{5}$=2}
\mcq{65}{The logarithmic form of $10^{3}=1000$ is:}{$\log_{10}3=1000$}{$\log_{3}10=1000$}{$\log_{10}1000=3$}{$\log_{3}1000=10$}
\mcq{66}{Solve $\log_{3}27=x$ :}{2}{9}{6}{3}
\mcq{67}{Convert 7$^{0}$=1 to logarithmic form:}{$\log_{1}7=0$}{$\log_{0}7=1$}{$\log_{7}1=0$}{$\log_{7}0=1$}
\mcq{68}{The base of $\log_{10}100=2$ :}{2}{10}{100}{1}
\mcq{69}{The characteristic of log 0.0078 is:}{-2}{-3}{3}{2}
\mcq{70}{The mantissa of log 345 (given log 345=2.5378) is:}{2}{0.5378}{5378}{2.5378}
\mcq{71}{If $\log x=\overline{3}.678$, the characteristic is:}{3}{-3}{-4}{4}
\mcq{72}{If log x = 2.3010, then x=:}{20}{200}{0.2}{2000}
\mcq{73}{The base of natural logarithm is:}{e}{10}{2}{3}
\mcq{74}{Simplify $\log(a)+\log(b)-\log(c)$ :}{$\log\!\left(\dfrac{ab}{c}\right)$}{$\log(a+b-c)$}{$\log\!\left(\dfrac{ac}{b}\right)$}{$\log\!\left(\dfrac{a+b}{b}\right)$}
\mcq{75}{Expand $\log(x^{2}y^{3})$ :}{6log}{logx + logy}{5log}{2logx + 3logy}
\mcq{76}{$\log_{3}\!\left(\dfrac{27}{9}\right)$ = ......:}{$\log_{3}27-\log_{3}9$}{$\log_{3}27+\log_{3}9$}{$\log_{3}(27\times 9)$}{$\log_{3}(27+3)$}
\mcq{77}{Simplify $\log_{5}2^{4}$ :}{$4\log_{5}2$}{$2\log_{5}4$}{$\log_{5}8$}{$\log_{5}16$}
\mcq{78}{$\log_{4}8+\log_{4}2=$ :}{$\log_{8}2$}{$\log_{4}10$}{$\log_{8}4$}{$\log_{4}16$}
\mcq{79}{Change $\log_{7}5$ :}{$\log_{7}10$}{$\dfrac{\log_{5}7}{\log_{5}10}$}{$\dfrac{\log_{10}5}{\log_{10}7}$}{$\log_{10}7$}
\mcq{80}{Simplify $\log_{6}6^{5}$ :}{5}{6}{1}{0}
\mcq{81}{$\log_{3}81-\log_{3}9=$}{$\log_{3}9$}{$\log_{3}72$}{$\log_{3}(81/9)$}{$\log_{3}(81\times 9)$}
\mcq{82}{If $\log_{a}b=c$ then}{$c^{a}=b$}{$b^{c}=a$}{$a^{b}=c$}{$a^{c}=b$}
\mcq{83}{Convert Ln(x)=2 to exponential form:}{e$^{2}$=x}{10$^{2}$=x}{2$^{10}$=x}{x$^{2}$=e}
\mcq{84}{If $M=\log(A/A_{o})$, If A=1000 and A$_{o}$=10, find M:}{4}{3}{2}{5}
\mcq{85}{The standard form of 5.2 x 10$^{6}$ is:}{52,000}{520,000}{5,200,000}{52,000,000}
\mcq{86}{Scientific notation of 0.00034 is:}{3.4 x 10$^{3}$}{3.4 x 10$^{-4}$}{3.4 x 10$^{4}$}{3.4 x 10$^{-3}$}
\mcq{87}{The base of common logarithm is:}{2}{10}{5}{e}
\mcq{88}{Log$_{2}$2$^{3}$:}{1}{2}{5}{3}
\mcq{89}{Log 100 = ......:}{2}{3}{10}{1}
\mcq{90}{If log 2 = 0.3010, then log 200 is:}{1.3010}{0.6010}{2.3010}{2.6010}
\mcq{91}{Log (0) = ......:}{Positive}{Negative}{Zero}{Undefined}
\mcq{92}{log 10,000 =.......:}{2}{3}{4}{5}
\mcq{93}{log 5 + log 3= ......:}{Log 0}{Log 2}{$\log\!\left(\dfrac{5}{3}\right)$}{Log 15}
\mcq{94}{3$^{4}$= 81 in logarithmic form is:}{log$_{3}$4=81}{log$_{4}$3 = 81}{log$_{3}$81 = 4}{log$_{4}$81 = 3}
\mcq{95}{Sequence 0,1,1,2,3,5,8,13,21, ......, example of ...... sequence:}{Roman}{Odd}{Even}{Fibonacci}
\mcq{96}{Georg cantor was ...... mathematician:}{Japani}{German}{Greek}{Egyptian}
\mcq{97}{Fibonacci sequence Formula F$_{n}$= ......:}{F$_{n-1}$+F$_{n+2}$}{F$_{n-1}$+F$_{n-2}$}{F$_{n-1}$-F$_{n+2}$}{F$_{n-1}$+F$_{n}$}
\mcq{98}{Name of sets represents:}{1,2,3,...}{A,B,C ...}{Both A and B}{None of these}
\mcq{99}{How methods to describe a set:}{1}{2}{3}{4}
\mcq{100}{SetA=\{a,e,i,o,u\} written in form:}{Descriptive}{Tabular}{Builder}{None of these}
\mcq{101}{Symbol use for member or belong of a set:}{$\subseteq$}{$\in$}{$\notin$}{$\forall$}
\mcq{102}{Set \{2,3,5,7,11,......\} is called ....... member:}{Natural}{whole}{Odd}{Prime}
\mcq{103}{Q$\cup$Q$'$=......:}{N}{R}{Q}{Q$'$}
\mcq{104}{The set\{0\} is:}{Singleton}{Empty}{Null}{Prime}
\mcq{105}{The set with no element is called ...... set:}{Singleton set}{Empty}{Power}{Prime}
\mcq{106}{Symbol used for empty set:}{$\{\phi\}$}{$\phi$}{\{\{ \}\}}{None of these}
\mcq{107}{Example of equal set:}{\{a,b,c\}, \{1,2,3\}}{\{1,2,3\}, \{2,1,3\}}{Both a and b}{None of these}
\mcq{108}{Symbol used for equivalent set A and B:}{A=B}{$A\equiv B$}{A+B}{A$\cup$B}

\mcq{109}{If A subset B them written as sympetaly:}{$A\subseteq B$}{$A\supseteq B$}{$A\in B$}{$A\notin B$}
\mcq{110}{Elements of power set C=\{a,b,c,d\}:}{8}{16}{24}{26}
\mcq{111}{If D=\{a\} then, P(D)=......:}{\{a\}}{$\phi,\{a\}$}{$\{\phi,\{a\}\}$}{None of these}
\mcq{112}{S$_{1}$= set of odd number, S$_{2}$= set of even number the intersection of both sets:}{Universal set}{Empty set}{Even set}{Odd set}
\mcq{113}{$\{x\mid x\in A\ \vee\ x\in B\}$ represent as:}{A$\cup$B}{$A\cap B$}{A-B}{$\phi$}
\mcq{114}{Symbol used for intersection:}{U}{I}{$\cap$}{-}
\mcq{115}{$A\cap B=......$ :}{$\{x\mid x\in A\vee x\in B\}$}{$\{x\mid x\in A\wedge x\in B\}$}{$\{x\mid x\in A\wedge x\notin B\}$}{$\{x\mid x\notin A\wedge x\in B\}$}
\mcq{116}{$A\cap B$=\{1,2,3\} if A=\{1,2,3,7\} then equation B=......}{\{1,2,6\}}{\{1,2,3,4\}}{\{1,3,6\}}{\{1,6,7\}}
\mcq{117}{If $A\cap B=\phi$ :}{Subset}{Overlapping}{Disjoint}{None of these}
\mcq{118}{Which example of overlapping set:}{\{1,2\}, \{2,3,4\}}{\{1,2\}, \{1,2,3,4\}}{\{1,2\}, \{3,4\}}{None of these}
\mcq{119}{$\{x\mid x\in A\vee x\notin B\}$ shows as:}{A-B}{B-A}{A$\cup$B}{A+B}
\mcq{120}{If A=\{1,2,3,4,5\} and B=\{4,5,6,7,8,9,10\} then, A-B = .......:}{\{4,5\}}{\{1,2,3\}}{\{7,8,9,10\}}{\{4,5,6,7,8,9,10\}}
\mcq{121}{A$'$......:}{U-A}{A-U}{U}{U+A}
\mcq{122}{A$'$=......:}{U-A}{A-U}{U}{$\phi$}
\mcq{123}{(A$\cup$B)$'$=......:}{$A\cap B$}{A$'\cup$B$'$}{$A'\cap B'$}{A$'\cup$B}
\mcq{124}{$A\cap B=B\cap A$ is called ...... property of intersection:}{Associative}{Commutative}{Distributive}{None of these}
\mcq{125}{If A=\{1,3,5,7,9,11\}, then n(A)= .......:}{6}{4}{5}{2}
\mcq{126}{$n(A)+n(B)-n(A\cap B)=......$}{$n(A\cap B)$}{$n(A\cup B)$}{n}{None of these}
\mcq{127}{A x B write symbolically as:}{$\{(x,y)\mid x\in A$ and $x\in B\}$}{$\{(x,y)\mid x\in A$ and $y\in B\}$}{$\{(x,y)\mid x\notin A$ and $x\in B\}$}{$\{(x,y)\mid x\in B$ and $y$}
\mcq{128}{Let A=\{1,2,3\}, determine the relation $r=\{(x,y\mid x<y)\}$ :}{\{(2,1),(3,1),(3,2)\}}{\{(1,2),(1,3),(2,3)\}}{\{(1,1),(2,2),(3,3)\}}{None of these}
\mcq{129}{If $r$=\{(1,2),(1,3),(2,3)\}, then $r$=:}{\{1,2\}}{\{1\}}{\{2,3\}}{\{1,2,3\}}
\mcq{130}{Each ordered pair consist of ...... coordinates:}{1}{2}{3}{4}
\mcq{131}{In ordered pair (x,y) x is called:}{Abscissa}{Ordinate}{Both}{None of these}
\mcq{132}{If a function $f:A\rightarrow B$ such that range $f\neq B$ the f is called ...... function.}{Onto}{Into}{One-one}{Objective}
\mcq{133}{If $f(n)=2n-1$ then $f(7)$ is:}{12}{13}{14}{15}
\mcq{134}{If $g(n)=x^{2}-3$ then g(4) is:}{7}{-7}{13}{9}
\mcq{135}{The set builder form of the set $\left\{1,\dfrac{1}{3},\dfrac{1}{5},\dfrac{1}{7},\ldots\right\}$ :}{$\left\{x\mid x=\dfrac{1}{n},n\in W\right\}$}{$\left\{x\mid x=\dfrac{1}{2n+1},n\in W\right\}$}{$\left\{x\mid x=\dfrac{1}{n+1},n\in W\right\}$}{$\{x\mid x=2n+1,\ n\in$}
\mcq{136}{If A = \{ \} , then P(A) is:}{\{ \}}{\{1\}}{\{\{ \}\}}{$\phi$}
\mcq{137}{If $U=\{1,2,3,4,5\}$, $A=\{1,2,3\}$ and $B=\{3,4,5\}$, then $U-(A\cap B)$ is:}{\{1,2,4,5\}}{\{2,3\}}{\{1,3,4,5\}}{\{1,2,3\}}
\mcq{138}{If A and B are overlapping sets, then $n(A-B)$ :}{n(A)}{n(B)}{$A\cap B$}{$n(A)-n(A\cap B)$}
\mcq{139}{If $A\subseteq B$ and $B-A\neq\phi$, then $n(B-A)$ is equal to:}{0}{-n(B)}{n(A)}{n(B)-n(A)}
\mcq{140}{If $n(A\cup B)=50$, $n(A)=30$ and $n(B)=35$, then $n(A\cap B)$=:}{23}{15}{9}{40}
\mcq{141}{If $A=\{1,2,3,4\}$ and $B=\{x,y,z\}$, then Cartesian product of A and B contains exactly ...... elements:}{13}{12}{10}{6}
\mcq{142}{If $f(x)=x^{2}-3x+2$, then the value of $f(a+1)$ is equal to:}{a+1}{a$^{2}$+1}{a$^{2}$+2a+1}{a$^{2}$-a}
\mcq{143}{Given that $f(x)=3x+1$, if $f(x)=28$, then the value of x is:}{9}{27}{3}{18}
\mcq{144}{Let $A=\{1,2,3\}$ and $B=\{a,b\}$ two non-empty sets and $f:A\rightarrow B$ be a function defined as $f=\{(1,a),(2,b),(3,b)\}$, then which of the following statement is true?}{f is injective}{f is surjective}{f is bijective}{f is into only}
\mcq{145}{Common factor of 2x - 6 is:}{2}{3}{x}{None of these}
\mcq{146}{Factors of x$^{2}$-3x-10 are:}{(x + 2)(x + 5)}{(x + 2)(x - 5)}{(x - 2)(x + 5)}{(x - 2)(x - 5)}
\mcq{147}{Which expression type of quadratic expression:}{ax + b}{ax$^{2}$+ bx + c}{ax$^{3}$+ bx$^{2}$+ c}{None of these}
\mcq{148}{Which pair is suitable for factorization of x$^{2}$-11x +24:}{+3, +8}{-3,-8}{-3, +8}{+3,-8}
\mcq{149}{An expression having degree 2 is called ...............expression:}{Cubic}{Linear}{Quadratic}{None of these}
\mcq{150}{Factors of 3x$^{2}$-4x-4 are:}{(3x + 2)(x - 2)}{(3x - 2)(x + 2)}{(3x + 2)(x + 2)}{None of these}
\mcq{151}{Which of the following is a prime number:}{9}{15}{13}{21}
\mcq{152}{x$^{2}$-2ab+b$^{2}$is equal to:}{(a + b)$^{2}$}{(a - b)$^{2}$}{a$^{2}$- b$^{2}$}{a$^{2}$+ b$^{2}$}
\mcq{153}{Factors of 8$^{3}$+27:}{(2x + 3)(4x$^{2}$+ 6x + 9)}{(2x + 3)(4x$^{2}$- 6x + 9)}{(2x - 3)(4x$^{2}$+ 6x + 9)}{(2x - 3)(4x$^{2}$- 6x + 9)}
\mcq{154}{In 2x$^{2}$+7x+3 one factor is 2x + 1 then other factor:}{(x + 2)}{(x + 1)}{(x + 3)}{(x + 4)}
\mcq{155}{Expression with two term is called:}{Monomial}{Binomial}{Trinomial}{None of these}
\mcq{156}{a$^{3}$+b$^{3}$+3ab(a+b) is equal to:}{(a + b)$^{3}$}{(a - b)$^{3}$}{a$^{3}$- b$^{3}$}{a$^{3}$+ b$^{3}$}
\mcq{157}{How many method to solve H.C.F:}{1}{2}{3}{4}
\mcq{158}{H.C.F of 6x$^{2}$y,9xy$^{2}$:}{3x$^{2}$y}{3xy}{6x$^{2}$y}{54x$^{2}$y$^{3}$}
\mcq{159}{L.C.M = ..................:}{$\dfrac{p(x)\times q(x)}{H.C.F}$}{$\dfrac{H.C.F\ \times\ p(x)}{q(x)}$}{P}{None of these}
\mcq{160}{Square root of a$^{2}$-2a+1:}{$\pm$(a-1)}{$\pm$(a+1)}{a - 1}{a + 1}
\mcq{161}{The square root of x$^{2}$-8x+16 is:}{(x+4)}{-(x+4)}{Both A and B}{None of these}
\mcq{162}{The return on investment R(x) is given by R(x) =x$^{2}$+6x-8 What are the investment levels that result is zero return:}{x=2, x=4}{x=1, x=8}{x=2, x=-4}{x=-2, x=4}
\mcq{163}{The factorization of 12x +36 is:}{12(x+3)}{12(3x)}{12(3x+1)}{x(12+36x)}
\mcq{164}{The factors of $4x^{2}-12y+9$ :}{$(2x+3)^{2}$}{$(2x-3)^{2}$}{(2x-3) (2x+3)}{$(2+3x)(2-3x)^{2}$}
\mcq{165}{The HCF of $a^{3}b^{3}$ :}{$a^{3}b^{3}$}{$ab^{2}$}{$a^{4}b^{5}$}{$a^{2}b$}
\mcq{166}{The LCM of $16x^{2}$, 4x and 30xy is:}{$480x^{3}y$}{240xy}{$240x^{2}y$}{$120x^{4}y$}
\mcq{167}{Product of LCM and HCF = ...... of two polynomials:}{Sum}{Difference}{Product}{Quotient}
\mcq{168}{The square root of $x^{2}-6x+9$ :}{$\pm(x-3)$}{$\pm(x+3)$}{x-3}{x+3}
\mcq{169}{The LCM of $(a-b)^{2}$ and $(a-b)^{4}$ is:}{$(a-b)^{2}$}{$(a-b)^{3}$}{$(a-b)^{4}$}{$(a-b)^{6}$}
\mcq{170}{Factorization of $x^{3}+3x^{2}+3x+1$ :}{$(x+1)^{3}$}{$(x-1)^{3}$}{$(x+1)(x^{2}+x+1)$}{$(x-1)(x^{2}-x+1)$}
\mcq{171}{Cubic polynomial has degree:}{1}{2}{3}{4}
\mcq{172}{One of the factors of $x^{3}-27$ :}{x-3}{x+3}{$x^{2}-3x+9$}{Both a and c}
\mcq{173}{Which of the following is a linear equation in one variable:}{2x + 3y = 7}{5x - 4 = 0}{x$^{2}$+ 3x = 8}{xy + 2 = 5}
\mcq{174}{The solution of the equation 3x - 7 = 2 is:}{3}{-3}{$\dfrac{5}{3}$}{9}
\mcq{175}{In the inequality 4x+2>10,what is the solution for x:}{x > 2}{x < 2}{x >1}{x < 2}
\mcq{176}{Which symbol represents 'less than or equal to'?4.\urdu{کم ہے اور برابر کو کونسی علامت ظاہر کرتی ہے؟}}{$>$}{$<$}{$\leq$}{$\geq$}
\mcq{177}{The graphical representation of a linear equation in two variables is:}{A curve}{A straight line}{A parabola}{A circle}
\mcq{178}{What is the solution of 2(x-3)=4:}{x = 2}{x = 5}{x = 3}{x = 1}
\mcq{179}{The inequality $x-5\geq 3$ means:}{x is greater than 3}{x is greater than or equal to 3}{x is greater than or equal to 8}{x is less than 3}
\mcq{180}{The equation x+2y=6 represents a:}{Quadratic equation}{Linear equation}{Exponential equation}{None of these}
\mcq{181}{The associated equation of the inequality $x+2y\leq 8$ :}{x + 2y = 8}{x + 2y $>$ 8}{x + 2y $<$ 8}{$x+2y\geq 8$}
\mcq{182}{The solution of the system x - y = 2 and x + y = 6 is:}{(4, 2)}{(2, 4)}{(3, 3)}{(6, 0)}
\mcq{183}{If x = 3 is a solution of 5x - 2 = k, then k is:}{11}{13}{15}{17}
\mcq{184}{What is the graphical representation of a system of linear inequalities:}{A shaded region}{A single point}{A straight line}{A curve}
\mcq{185}{The equation 2x - y = 4 is equivalent to the inequality:}{2x - y $>$ 4}{2x - y $<=$ 4}{2x - y $>=$ 4}{None of these}
\mcq{186}{Which of the following is an objective function:}{x + y}{x$^{2}$ + y$^{2}$}{3x + 2y}{Both a and c}
\mcq{187}{The feasible region for an inequality system is always:}{A single point}{A region on the plane}{An empty set}{A line segment}
\mcq{188}{A vertical line divided the plane into:}{Two half-planes}{Four quadrants}{Three sections}{Infinite parts}
\mcq{189}{If a system of inequalities has no solution, its feasible region is:}{Infinite}{A single point}{A line segment}{Empty}
\mcq{190}{The intersection of two half-planes is known as:}{A vertex}{A feasible region}{A line}{A parabola}

\mcq{191}{The graph of $x+y\geq 5$ includes the line:}{x + y = 5}{x + y $<$ 5}{x + y $>$ 5}{None of these}
\mcq{192}{Which method is used to find the optimal solution of a linear programming problem:}{Graphical method}{Algebraic method}{Simplex method}{All of these}
\mcq{193}{The equation x = 2 represents a:}{Horizontal line}{Vertical line}{Diagonal line}{None of these}
\mcq{194}{A point that satisfies all given inequalities is called:}{Feasible solution}{Optimal solution}{Boundary point}{None of these}
\mcq{195}{The line y = 3x + 2 has a slope of:}{2}{3}{-3}{-2}
\mcq{196}{If the inequalities $x\geq 0$ and $y\geq 0$ are given, the solution region lies in:}{First quadrant}{Second quadrant}{Third quadrant}{Fourth quadrant}
\mcq{197}{The inequality 5x - 3y $<$ 15 represents:}{A parabola}{A straight line}{A shaded region}{A single point}
\mcq{198}{The maximum or minimum value of a function occurs at:}{A random point}{A feasible point}{A corner point}{An undefined point}
\mcq{199}{If $x-y\leq 4$ then the region of solution is:}{Above the line}{Below the line}{On the line}{None of these}
\mcq{200}{The intersection of x + y = 4 and x - y = 2 is:}{(1,3)}{(3,1)}{(2,2)}{(4,0)}
\mcq{201}{In the following, linear equation is:}{5x $>$ 7}{4x - 2 $<$ 1}{2x + 1 = 1}{4 = 1 + 3}
\mcq{202}{Solution of 5x - 10 = 10 is:}{0}{50}{4}{-4}
\mcq{203}{If 7x + 4 $<$ 6x + 6, then x belongs to the interval:}{$(2,\infty)$}{$[2,\infty)$}{$(-\infty,2)$}{$(-\infty,2]$}
\mcq{204}{A vertical line divides the plane into:}{Left half plane}{Right half plane}{Full plane}{Two half planes}
\mcq{205}{The linear equation formed out of the linear inequality is called:}{Linear equation}{Associated equation}{Quadratic equal}{None of these}
\mcq{206}{3x + 4 $<$ 0 is:}{Equation}{Inequality}{Not inequality}{Identity}
\mcq{207}{Corner point is also called:}{Code}{Vertex}{Curve}{Region}
\mcq{208}{(0,0) is solution of inequality:}{4x + 5y $>$ 8}{3x + y $>$ 6}{-2x + 3y $<$ 0}{x + y $>$ 4}
\mcq{209}{The solution region restricted to the first quadrant is called:}{Objective region}{Feasible region}{Solution region}{Constraints region}
\mcq{210}{A function that is to be maximized or minimized is called:}{Solution function}{Objective function}{Feasible function}{None of these}
\mcq{211}{In $\angle AOB$ vertex is:}{A}{O}{B}{None of these}
\mcq{212}{ray OA written as symbolically.}{$\overleftrightarrow{OA}$}{$\widehat{OA}$}{$\overline{OA}$}{$\overrightarrow{OA}$}
\mcq{213}{Co-terminal of 65$^{\circ}$:}{360$^{\circ}$}{-360$^{\circ}$}{425$^{\circ}$}{-405$^{\circ}$}
\mcq{214}{One full revolution of a circle is ............radian:}{$\pi$}{$2\pi$}{$3\pi$}{$4\pi$}
\mcq{215}{Area of sector is = ................:}{$\dfrac{1}{2}r^{2}\theta$}{$\pi r^{2}$}{$4\pi r^{2}$}{$\dfrac{1}{2}r^{3}\theta$}
\mcq{216}{Value of $\dfrac{\pi}{4}$ radian:}{$\sqrt{2}$}{$\sqrt{3}$}{$\dfrac{1}{\sqrt{2}}$}{$\dfrac{1}{\sqrt{3}}$}
\mcq{217}{What is the measure of a full rotation in degrees:}{90$^{\circ}$}{180$^{\circ}$}{270$^{\circ}$}{360$^{\circ}$}
\mcq{218}{What is the measure of a right angle in radians?}{$\dfrac{\pi}{2}$}{$\pi$}{$\dfrac{3\pi}{2}$}{$2\pi$}
\mcq{219}{Which of the following is an obtuse angle:}{45$^{\circ}$}{90$^{\circ}$}{120$^{\circ}$}{180$^{\circ}$}
\mcq{220}{What is the measure of a straight angle in radians?}{$2\pi$}{$\dfrac{3\pi}{2}$}{$\pi$}{$\dfrac{\pi}{2}$}
\mcq{221}{What is the measure of 1 radian in degrees:}{57.2958$^{\circ}$}{90$^{\circ}$}{180$^{\circ}$}{360$^{\circ}$}
\mcq{222}{Convert 73.12$^{\circ}$to degrees, minutes, and seconds:}{73$^{\circ}$12$'$ 12$''$}{73$^{\circ}$12$'$ 7$''$}{73$^{\circ}$7$'$ 7$''$}{73$^{\circ}$7$'$ 12$''$}
\mcq{223}{Convert 45$^{\circ}$ 45$'$ 45$''$to decimal degrees:}{45.1234$^{\circ}$}{45.4567$^{\circ}$}{45.7625$^{\circ}$}{45.9876$^{\circ}$}
\mcq{224}{What is the value of 1$^{\circ}$in seconds:}{60$''$}{360$''$}{3600$''$}{60$'$}
\mcq{225}{What is the measure of a full circle in radians?}{$2\pi$}{$\dfrac{3\pi}{2}$}{$\pi$}{$\dfrac{\pi}{2}$}
\mcq{226}{Convert $\dfrac{5\pi}{3}$ radians to degrees.}{300$^{\circ}$}{150$^{\circ}$}{210$^{\circ}$}{330$^{\circ}$}
\mcq{227}{Convert 15$^{\circ}$to radians:}{$\dfrac{\pi}{12}$}{$\dfrac{\pi}{6}$}{$\dfrac{\pi}{4}$}{$\dfrac{\pi}{3}$}
\mcq{228}{Convert 1.2 radians to degrees:}{90$^{\circ}$}{57.2958$^{\circ}$}{45$^{\circ}$}{68.75$^{\circ}$}
\mcq{229}{In right angled triangle ABC if $m\angle B=90^{o}$ then which statement is true:}{a$^{2}$+ b$^{2}$= c$^{2}$}{a$^{2}$- b$^{2}$= c$^{2}$}{a$^{2}$+ b$^{2}$= b$^{2}$}{a$^{2}$+ b$^{2}$=a$^{2}$}
\mcq{230}{Reciprocal of $\sin\theta$ :}{$\sec\theta$}{$\mathrm{cosec}\,\theta$}{$\cot\theta$}{$\tan\theta$}
\mcq{231}{cot $\theta$ =...................}{$\dfrac{1}{\sin\theta}$}{$\dfrac{1}{\cos\theta}$}{$\dfrac{1}{\tan\theta}$}{1-cos$^{2}\theta$}
\mcq{232}{sin $\theta$ =................... :}{$\sqrt{1-\cos^{2}\theta}$}{$\dfrac{1}{\tan\theta}$}{$\dfrac{\sin\theta}{\cos\theta}$}{$1-\cos^{2}\theta$}
\mcq{233}{If sin $\theta=\dfrac{2}{3}$ then value of cosec $\theta$ :}{$\dfrac{2}{3}$}{$\dfrac{-2}{3}$}{$\dfrac{3}{2}$}{$\dfrac{-3}{2}$}
\mcq{234}{Fundamental trigonometric ratio are:}{3}{4}{5}{6}
\mcq{235}{What is the value of cot 30$^{\circ}$ :}{$\sqrt{3}$}{$\dfrac{1}{\sqrt{3}}$}{1}{0}
\mcq{236}{What is the value of sin (90$^{\circ}$- $\theta$)?}{Tan$\theta$}{Sin$\theta$}{Cos$\theta$}{Cot$\theta$}
\mcq{237}{What is the value of tan $\theta$ in terms of sin $\theta$ and cos $\theta$?}{$\dfrac{\sin\theta}{\cos\theta}$}{$\dfrac{\cos\theta}{\sin\theta}$}{$\dfrac{1}{\sin\theta}$}{$\dfrac{1}{\cos\theta}$}
\mcq{238}{Which following is a Pythagorean identity:}{$\sin^{2}\theta+\cos^{2}\theta=1$}{$\sin\theta+\cos\theta=1$}{$\tan\theta+\cot\theta=1$}{$\sec\theta+\mathrm{cosec}\,\theta=1$}
\mcq{239}{What is the value of tan$^{2}$ $\theta$ + 1?}{Sin$^{2}\theta$}{Cosec$^{2}\theta$}{Cot$^{2}\theta$}{Sec$^{2}\theta$}
\mcq{240}{What is the value of 1 + cot$^{2}\theta$ :}{Cosec$^{2}\theta$}{Sec$^{2}\theta$}{Tan$^{2}\theta$}{Sin$^{2}\theta$}
\mcq{241}{If $\tan\theta=\dfrac{3}{4}$, what is the value of $\sin\theta$}{$\dfrac{5}{3}$}{$\dfrac{4}{5}$}{$\dfrac{3}{5}$}{$\dfrac{5}{4}$}
\mcq{242}{What is the value of $\sin 45^{\circ}$?}{$\dfrac{1}{\sqrt{2}}$}{$\dfrac{\sqrt{3}}{2}$}{$\dfrac{1}{2}$}{1}
\mcq{243}{What is the value of cos 30$^{\circ}$ :}{$\dfrac{1}{2}$}{$\dfrac{\sqrt{3}}{2}$}{$\dfrac{1}{\sqrt{2}}$}{1}
\mcq{244}{What is the value of tan 60$^{\circ}$?}{$\sqrt{3}$}{$\dfrac{1}{\sqrt{3}}$}{1}{$\dfrac{\sqrt{3}}{2}$}
\mcq{245}{What is the value of cot 45$^{\circ}$ :}{1}{$\dfrac{\sqrt{3}}{2}$}{$\dfrac{1}{\sqrt{3}}$}{0}
\mcq{246}{In a right-angled triangle, if one angle is 30$^{\circ}$ what is the other acute angle:}{60$^{\circ}$}{45$^{\circ}$}{90$^{\circ}$}{120$^{\circ}$}
\mcq{247}{In a right-angled triangle, if one side is 3 cm and the other side is 4 cm, what is the length of the hypotenuse:}{5 cm}{6 cm}{7 cm}{8 cm}
\mcq{248}{In a right-angled triangle, if one angle is 45$^{\circ}$ what is the length of the hypotenuse if one side is 5 cm?}{$5\sqrt{2}$ cm}{5 cm}{10 cm}{$10\sqrt{2}$}
\mcq{249}{In a right-angled triangle, if the hypotenuse is 10 cm and one side is 6 cm, what is the length of the other side:}{4 cm}{6 cm}{8 cm}{2 cm}
\mcq{250}{In a right-angled triangle, if one side is 3 cm and the other side is 4 cm, what is the length of the hypotenuse:}{5 cm}{6 cm}{7 cm}{8 cm}
\mcq{251}{The value of $\tan^{-1}2$ :}{$\dfrac{\pi}{2}$}{$\dfrac{3\pi}{2}$}{0.4636 $\pi$}{0.4636}
\mcq{252}{In a right triangle, the hypotenuse is 13 units and one of the angles is $\theta=30^{\circ}$ :}{6.5 units}{7.5 units}{6 units}{5 units}
\mcq{253}{A person standing 50 m away from a building sees the top of the building at an angle of elevation of 45$^{\circ}$ Height of the building is:}{50 m}{25 m}{35 m}{70 m}
\mcq{254}{$\sec^{2}\theta-\tan^{2}\theta$ :}{$\sin^{2}\theta$}{1}{$\cos^{2}\theta$}{$\cot^{2}\theta$}
\mcq{255}{If $\sin\theta=\dfrac{3}{5}$ and $\theta$ is an acute angle, $\cos^{2}\theta$= ......:}{$\dfrac{7}{25}$}{$\dfrac{24}{25}$}{$\dfrac{16}{25}$}{$\dfrac{4}{25}$}
\mcq{256}{$\dfrac{5\pi}{24}$ rad =.........degree}{30$^{\circ}$}{37.5$^{\circ}$}{45$^{\circ}$}{52.5$^{\circ}$}
\mcq{257}{292.5$^{\circ}$= ......rad:}{$\dfrac{17\pi}{6}$}{$\dfrac{17\pi}{4}$}{1.6 $\pi$}{1.625 $\pi$}
\mcq{258}{Which of the following is a valid identity:}{$\cos\!\left(\dfrac{\pi}{2}-\theta\right)=\sin\theta$}{$\cos\!\left(\dfrac{\pi}{2}-\theta\right)=\cos\theta$}{$\cos\!\left(\dfrac{\pi}{2}-\theta\right)=\sec\theta$}{$\cos\!\left(\dfrac{\pi}{2}-\theta\right)=\cos$}
\mcq{259}{Sin 60$^{\circ}$= ......:}{1}{$\dfrac{1}{2}$}{$\sqrt{(3)^{2}}$}{$\dfrac{\sqrt{3}}{2}$}
\mcq{260}{$\cos^{2}100\,\pi+\sin^{2}100\,\pi$ :}{1}{2}{3}{4}
\mcq{261}{In a plane horizontal line represent:}{xox$'$}{yoy$'$}{zoz$'$}{None of these}
\mcq{262}{In a plane yoy$'$ is called .............:}{Horizontal line}{Vertical line}{Parallel line}{None of these}
\mcq{263}{The points laying on Ox is:}{+ ve}{- ve}{+ ve or - ve}{None of these}
\mcq{264}{In (x, y) if x $>$ 0 and y $<$ 0 then point in which quadrant line:}{I}{II}{III}{IV}
\mcq{265}{The point P(3,- 4) lies in which quadrant:}{I}{II}{III}{IV}

\mcq{266}{The distance between the points A(2, 3) and B(5, 7) Is:}{5}{6}{7}{8}
\mcq{267}{The midpoint of the line segment joining (1, 2) and (5, 8) is:}{(3, 5)}{(4, 5)}{(3, 4)}{(4,6)}
\mcq{268}{If the distance between (x, 3) and (4, 7) is 5, then x is:}{1}{2}{3}{4}
\mcq{269}{The point (0, 5) lies on which axis:}{x-axis}{y-axis}{Both axes}{None}
\mcq{270}{The midpoint of a bridge connecting (2, 5) and (8, 9) is:}{(5,7)}{(6,7)}{(5,8)}{(6,8)}
\mcq{271}{The slope of the line passing through (2, 3) and (4, 7) is:}{1}{2}{3}{4}
\mcq{272}{The inclination of a line with slope 1 is:}{30$^{\circ}$}{45$^{\circ}$}{60$^{\circ}$}{90$^{\circ}$}
\mcq{273}{If the slope of a line is undefined, the line is:}{Horizontal}{Vertical}{Parallel to x-axis}{Parallel to y-axis}
\mcq{274}{The slope of a line parallel to the x-axis is:}{0}{1}{Undefined}{-1}
\mcq{275}{The slope of a line perpendicular to a line with slope 2 is:}{2}{-2}{$\dfrac{1}{2}$}{$-\dfrac{1}{2}$}
\mcq{276}{The equation of a line with slope 3 and y-intercept 2 is:}{y = 3x + 2}{y = 2x + 3}{y = 3x - 2}{y = 2x - 3}
\mcq{277}{The equation of a line passing through (1, 2) with slope 4 is:}{y = 4x - 2}{y = 4x + 2}{y = 4x - 6}{y = 4x + 6}
\mcq{278}{The equation of a line passing through (2, 3) and (4, 7) is:}{y = 2x - 1}{y = 2x + 1}{y = 3x - 1}{y = 3x + 1}
\mcq{279}{The equation of a line with x-intercept 3 and y-intercept 4 is:}{4x+3y=12}{3x+4y=12}{4x-3y=12}{3x- 4y=12}
\mcq{280}{The equation of a line parallel to y = 2x + 3 and passing through (1, 1) is:}{y = 2x - 1}{y = 2x + 1}{y = 2x - 3}{y = 2x + 3}
\mcq{281}{The equation of a horizontal line passing through (3, 4) is:}{y = 3}{y = 4}{x = 3}{x = 4}
\mcq{282}{The equation of a vertical line passing through (5, 6) is:}{y = 5}{y = 6}{x = 5}{x = 6}
\mcq{283}{The slope of the line 2x + 3y = 6 is:}{$\dfrac{2}{3}$}{$-\dfrac{2}{3}$}{$\dfrac{3}{2}$}{$-\dfrac{3}{2}$}
\mcq{284}{If two lines are perpendicular, the product of their slopes is:}{1}{-1}{0}{2}
\mcq{285}{The slope of a line perpendicular to x = 3 is:}{0}{1}{2}{-1}
\mcq{286}{The equation 2x + 3y = 6 in slope-intercept form is:}{$y=\dfrac{-2}{3x+2}$}{$y=\dfrac{2}{3x+2}$}{$y=\dfrac{-2}{3x-2}$}{$y=\dfrac{2}{3x-2}$}
\mcq{287}{The equation 3x - 4y = 12 in intercept form is:}{$\dfrac{x}{4}+\dfrac{y}{3}=1$}{$\dfrac{x}{4}-\dfrac{y}{3}=1$}{$\dfrac{x}{3}+\dfrac{y}{4}=1$}{$\dfrac{x}{3}-\dfrac{y}{4}=1$}
\mcq{288}{The equation 5x - 12y + 39 = 0 in normal form is:}{5x - 12y = - 39}{5x + 12y = 39}{5x - 12y = 39}{5x + 12y = - 39}
\mcq{289}{The distance between two towns with coordinates (3, 4) and (7, 1) is:}{5}{6}{7}{8}
\mcq{290}{The length of the side of a triangle with vertices A(2, 3) B(5, 7) and C(6,2) is:}{5}{6}{7}{8}
\mcq{291}{The distance between two cities with coordinates (10, 20) and (30, 40) is:}{$20\sqrt{2}$}{$30\sqrt{2}$}{$40\sqrt{2}$}{$50\sqrt{2}$}
\mcq{292}{The distance between two points on a map with coordinates (12, 5) and (8, -4) is:}{9}{10}{11}{12}
\mcq{293}{The length of a fence around a rectangular field with corners at (0,0) (0,5), (8,5) and (8,0) is:}{40}{30}{36}{26}
\mcq{294}{The mid point of a journey from (40$^{\circ}$N, 100$^{\circ}$W) to (50$^{\circ}$N, 80$^{\circ}$W) is:}{(50$^{\circ}$N, 90$^{\circ}$W)}{(45$^{\circ}$N, 80$^{\circ}$W)}{(45$^{\circ}$N, 90$^{\circ}$W)}{(40$^{\circ}$N, 90$^{\circ}$W)}
\mcq{295}{The equation of a straight line in the slope-intercept form is written as:}{y=m(x+c)}{y-y$_{1}$=m(x-x$_{1}$)}{y = c + mx}{ax + by + c = 0}
\mcq{296}{The gradients of two parallel lines are:}{Equal}{Zero}{Negative reciprocals of each other}{Always undefined}
\mcq{297}{If the product of the gradients of two lines is -1, then the lines are:}{Parallel}{Perpendicular}{Collinear}{Coincident}
\mcq{298}{Distance between two points P(1,2) and Q(4,6) is:}{5}{6}{$\sqrt{13}$}{4}
\mcq{299}{The midpoint of a line segment with endpoints (-2,4) and (6,-2) is:}{(4,2)}{(2,1)}{(1,1)}{(0,0)}
\mcq{300}{A line passing through points (1,2) and (4,5) is:}{y = x + 1}{y = 2x + 3}{y = 3x - 2}{y = x + 2}
\mcq{301}{The equation of a line in point-slope form is:}{y=m(x+c)}{y-y$_{1}$=m(x-x$_{1}$)}{y=c+mx}{ax+by+c=0}
\mcq{302}{2x+3y-6=0 in the slope-intercept form is:}{$y=\dfrac{-2}{3}x+2$}{$y=\dfrac{2}{3}x-2$}{$y=\dfrac{2}{3}x+1$}{$y=\dfrac{-2}{3}x-2$}
\mcq{303}{The equation of a line in symmetric form is:}{$\dfrac{x}{a}+\dfrac{y}{b}=1$}{$\dfrac{x-x_{1}}{1}+\dfrac{y-y_{1}}{m}=\dfrac{z-z_{1}}{1}$}{$ax+by+c=0$}{$y-y_{1}=m(x-x_{1})$}
\mcq{304}{The equation of a line in normal form is:}{y = mx + c}{$\dfrac{x}{a}+\dfrac{y}{b}=1$}{$\dfrac{x-x_{1}}{\cos\alpha}=\dfrac{y-y_{1}}{\sin\alpha}$}{$x\cos\alpha+y\sin\alpha=p$}
\mcq{305}{Which of the following expression is often related to inductive reasoning:}{Based on repeated experiments}{If and only if statements}{Statement if proven by a theorem}{Based on general principles}
\mcq{306}{Which of the following sentence describe deductive reasoning:}{General conclusion from a limited number of observations}{Based on repeated experiments}{Based on units of information that are accurate}{Draw conclusion from wellknown facts}
\mcq{307}{Which one the following statement is true:}{The set of integers if finite}{The sum of the interior angles of any quadrilateral is always 180$^{\circ}$}{$\dfrac{22}{7}\notin Q'$}{All isosceles triangles are equilateral triangles}
\mcq{308}{Which of the following statement is best to represent the negation of the statement "The stove is burning":}{The stove is not burning}{The stove is dim}{The stove is turned to low heat}{It is both burning and not burning}
\mcq{309}{The conjunction of two statements p and q is true when:}{Both p and q are false}{Both p and q are true}{Onle q is true}{Only p is true}
\mcq{310}{A conditional is regarded as false only when:}{Antecedent is true and consequent is false}{Consequent is true and antecedent is false}{Antecedent is true only}{Consequent is false only}
\mcq{311}{Contrapositive of $q\rightarrow p$ :}{$q\rightarrow\sim p$}{$\sim q\rightarrow p$}{$\sim p\rightarrow\sim q$}{$\sim q\rightarrow\sim p$}
\mcq{312}{The statement "Every integer greater than 2 is a sum of two prime number" is:}{Theorem}{Conjecture}{Axiom}{Postulates}
\mcq{313}{The statement " A straight line can be drawn between any two points" is:}{Theorem}{Conjecture}{Axiom}{Logic}
\mcq{314}{The statement " The sum of the interior angle of a triangle is 180$^{\circ}$" is:}{Converse}{Theorem}{Axiom}{Conditional}
\mcq{315}{A sentence or mathematical expression which may be true of false but not both is called:}{Argument}{Statement}{Proof}{Check}
\mcq{316}{The sum of measures of interior angles of a triangle:}{90$^{\circ}$}{180$^{\circ}$}{270$^{\circ}$}{360$^{\circ}$}
\mcq{317}{The circumference of a circle with radius r is:}{$2\pi r$}{$\pi r^{2}$}{$4\pi r^{2}$}{None of these}
\mcq{318}{Symbol use for 'or':}{$\wedge$}{$\vee$}{$\sim$}{$\equiv$}
\mcq{319}{Symbol stand for $\wedge$ :}{And \urdu{اور}}{Or \urdu{یا}}{Not \urdu{نہیں}}{Is equivalent \urdu{برابر ہے}}
\mcq{320}{If q is negation is written as:}{$\sim q$}{$\wedge q$}{$\vee q$}{$\Rightarrow q$}
\mcq{321}{The conjunction of two statements p and q is symbolically written as:}{$p\rightarrow\sim q$}{$p\vee q$}{$p\wedge q$}{$p\rightarrow q$}
\mcq{322}{Example of conjunction is:}{3$<$5 $\wedge$ 7$<$8}{3$<$3 $\wedge$ 10$<$10}{2$<$5 $\wedge$ 9$<$6}{None of these}
\mcq{323}{Triangle have two right angles or Lahoreis capital of Sindh given statement are called:}{Conjunction}{Disjunction}{Negation}{None of these}
\mcq{324}{A compound statement of the form $p\rightarrow q$ :}{Conjunction}{Disjunction}{Negation}{Conditional}
\mcq{325}{In $p\rightarrow q$ "p" is called:}{Hypothesis}{Conclusion}{Both on and}{None of these}
\mcq{326}{A conditional is regarded as ............... when antecedent is true and consequent is false:}{True}{False}{Negation}{None of these}
\mcq{327}{The statement $p\rightarrow q$ is called:}{Conditional}{Biconelitional}{Conjunction}{Disjunction}
\mcq{328}{$p\rightarrow q$ is given conditional then converse:}{$q\rightarrow p$}{$\sim p\rightarrow\sim q$}{$\sim q\rightarrow\sim p$}{None of these}
\mcq{329}{$p\rightarrow q$ is given conditional then $\sim q\rightarrow\sim p$ is called:}{Inverse}{Converse}{Contpositive}{None of these}
\mcq{330}{A statement which is true all possible values of the variable involved is lenoun:}{Tautology}{Absurdity}{Contingency}{None of these}
\mcq{331}{A statement which is False all possible values:}{Tautology}{Absurdity}{Contingency}{None of these}
\mcq{332}{$\dfrac{a}{b}+\dfrac{c}{d}=\ldots\ldots$}{$\dfrac{ac}{bd}$}{$\dfrac{a+c}{b+d}$}{$\dfrac{ad+bc}{bd}$}{$\dfrac{ab+cd}{bd}$}
\mcq{333}{(x-1)$^{2}$+ (x+1)$^{2}$= ...................:}{4(x$^{2}$+1)}{3(x$^{2}$+1)}{2(x$^{2}$+1)}{(x$^{2}$+1)}
\mcq{334}{There are more than sided closed figure is called:}{Circle}{Rectangle}{Triangle}{Polygon}
\mcq{335}{.................. figures have same shape but different size:}{Congruent}{Similar}{Equal}{None of these}
\mcq{336}{Example of similar ...................:}{Scalene triangle}{Acute angle triangle}{Equilateral triangle}{Obtuse angle triangle}
\mcq{337}{Symbol used for Similarity:}{$\equiv$}{$\sim$}{=}{$\neq$}
\mcq{338}{ABC and DEF are similar than written as symbolically:}{$\triangle ABC\equiv\triangle DEF$}{$\triangle ABC\sim\triangle DEF$}{$\triangle ABC>\triangle DEF$}{$\triangle ABC<\triangle DEF$}

\mcq{339}{Area of parallelogram:}{Base + altitude}{Base x altitude}{$\dfrac{1}{2}$(Base x altitude)}{Length x width}
\mcq{340}{What is the condition for two polygons to be similar:}{Corresponding angles are equal and Corresponding sides are equal}{Corresponding angles are equal and Corresponding sides are proportional}{Corresponding angles are proportional and Corresponding sides are equal}{Corresponding angles are proportional and Corresponding sides are proportional}
\mcq{341}{If the ratio of Corresponding sides of two triangles is 3:4, what is the ratio of their areas:}{3:4}{9:8}{9:16}{12:16}
\mcq{342}{In the triangles ABC and DEF, if AB=6 cm, DE= 9 cm and AC= 8 cm, what is DF if the triangles are similar:}{10 cm}{12 cm}{14 cm}{16 cm}
\mcq{343}{If base =8cm and altitude =4cm then area of parallelogram:}{12 cm$^{2}$}{4 cm$^{2}$}{18 cm$^{2}$}{32 cm$^{2}$}
\mcq{344}{If the area of a smaller polygon is 25 cm$^{2}$ and the ratio of Corresponding sides is 1:2, what is the area of the larger polygon:}{100 cm$^{2}$}{75 cm$^{2}$}{50 cm$^{2}$}{125 cm$^{2}$}
\mcq{345}{Formula use for exterior angle of polygon:}{$\dfrac{360^{o}}{n}$}{$\dfrac{180^{o}}{n}$}{$\dfrac{(n-2)\times 180^{o}}{n}$}{None of these}
\mcq{346}{Size of each interior angle of regular polygon interior angle:}{$\dfrac{(n-2)\times 180^{o}}{n}$}{$\dfrac{180^{o}}{n}$}{$\dfrac{360^{o}}{n}$}{None of these}
\mcq{347}{What is the sum of interior angles of a hexagon:}{720$^{\circ}$}{540$^{\circ}$}{900$^{\circ}$}{1080$^{\circ}$}
\mcq{348}{What is the measure of each interior angle of a regular pentagon:}{72$^{\circ}$}{108$^{\circ}$}{120$^{\circ}$}{144$^{\circ}$}
\mcq{349}{What is the measure of each exterior angle of a regular octagon:}{30$^{\circ}$}{90$^{\circ}$}{60$^{\circ}$}{45$^{\circ}$}
\mcq{350}{How many diagonals does a regular heptagon have:}{7}{21}{14}{28}
\mcq{351}{What is the sum of exterior angle of any polygon:}{180$^{\circ}$}{360$^{\circ}$}{540$^{\circ}$}{720$^{\circ}$}
\mcq{352}{In an equilateral triangle, what is the measure of each angle:}{45$^{\circ}$}{60$^{\circ}$}{90$^{\circ}$}{120$^{\circ}$}
\mcq{353}{In an isosceles triangle, if ne angle is 80$^{\circ}$, what is the measure of the their two angles:}{50$^{\circ}$}{60$^{\circ}$}{70$^{\circ}$}{80$^{\circ}$}
\mcq{354}{In a parallelogram, what is the relationship between opposite sides:}{Equal and parallel}{Equal but notparallel}{Parallel but not equal}{Neither equal nr parallel}
\mcq{355}{What is the measure of each angle in a rectangle:}{60$^{\circ}$}{120$^{\circ}$}{90$^{\circ}$}{180$^{\circ}$}
\mcq{356}{In a rhombus, what is the relationship between the diagonals:}{Equal}{Unequal}{Parallel}{Bisect each their at right angles}
\mcq{357}{If two polygons are similar then:}{Their corresponding angles are equal}{Their areas are equal}{Their volumes are equal}{Their corresponding sides are equal}
\mcq{358}{The ratio of the areas of two similar polygons is:}{Equal to the ratio of their perimeters}{Equal to the square of the ratio of their corresponding sides}{Equal to the cube of the ratio of their corresponding sides}{Equal to the sum of their corresponding sides}
\mcq{359}{If the volumes of two similar solids is $125\,cm^{3}$ and $27\,cm^{3}$, the ratio of their corresponding heights is ......:}{3:5}{5:3}{25:9}{9:25}
\mcq{360}{The exterior angle of regular pentagon is:}{40$^{\circ}$}{45$^{\circ}$}{60$^{\circ}$}{72$^{\circ}$}
\mcq{361}{A parallelogram has an area of 64 cm$^{2}$ and a similar parallelogram has an area of 144 cm$^{2}$, If a side of the smaller parallelogram is 8 cm, the corresponding side of the larger parallelogram is:}{10 cm}{12 cm}{18 cm}{16 cm}
\mcq{362}{The total number of diagonals in a polygon with 9 sides is:}{18}{21}{25}{27}
\mcq{363}{Two shapes are similar, and their radii are in the ratio 4;5. If the surface area of the larger sphere is $500\pi$ cm$^{2}$:}{$256\pi$ cm$^{2}$}{$320\pi$ cm$^{2}$}{$400\pi$ cm$^{2}$}{$405\pi$ cm$^{2}$}
\mcq{364}{A regular polygon has an exterior angle of 30$^{\circ}$. How many diagonals does the polygon have:}{54}{90}{72}{108}
\mcq{365}{In a regular hexagon, the ratio of the length of a diagonal to the side length is:}{$\sqrt{3}$ : 1}{2:1}{3:2}{2:3}
\mcq{366}{A regular polygon has an interior angle of 165$^{\circ}$. How many sides does it have:}{15}{16}{20}{24}
\mcq{367}{What is the general form of a linear function?}{$f(x)=mx+c$}{$f(x)=ax^{2}+bx+c$}{$f(x)=ax^{3}+bx^{2}+cx+d$}{$f(x)=ke^{2x}$}
\mcq{368}{In the linear function y=3x-5, what is the slop 'm':}{5}{3}{-3}{-5}
\mcq{369}{What is the y-intercept of the linear function y= -2x+8:}{-2}{2}{-8}{8}
\mcq{370}{The graph of y=x$^{2}$ represents:}{A straight line}{A parabola opening upwards}{A parabola opening downwards}{A cubic curve}
\mcq{371}{In the quadratic function y=-x$^{3}$, what is the direction of the parabola:}{Upwards}{Downwards}{Left}{Right}
\mcq{372}{The standard form of cubic functions:}{y=ax$^{3}$+bx$^{2}$+cx+d}{y=ax$^{2}$+bx+c}{y=ax+b}{None of these}
\mcq{373}{The graph of y=2x$^{3}$-3x$^{2}$+x-5 passes through which point:}{(0,-5)}{(0,0)}{(0,5)}{(0,-10)}
\mcq{374}{What is the general form of a reciprocal function?}{y=ax+b}{y=ax$^{2}$+bx+c}{$y=\dfrac{q}{x}$}{y=x$^{3}$}
\mcq{375}{What type of function is (equation):}{Exponential}{Linear}{Quadratic}{Cubic}
\mcq{376}{What does the graph of y=x$^{3}$ look like:}{A straight line}{A parabola}{A cubic curve passing through the origin}{A hyperbola}
\mcq{377}{The graph of the function y=5x represents:}{Exponential growth}{Exponential decay}{A line passing through the origin}{A hyperbola}
\mcq{378}{For the quadratic function y = x$^{2}$+ 2x - 3 what is the vertex:}{(-1, -3)}{(0, 3)}{(1, 2)}{(-2, 1)}
\mcq{379}{The graph of y = - x$^{2}$+ 4 cuts the y-axis at:}{(0, 4)}{(0, -4)}{(4, 0)}{(-4, 0)}
\mcq{380}{What is the slope of the line y = 2x + 3:}{2}{3}{-2}{0}
\mcq{381}{What is the slope of the line y = - 3x + 2:}{3}{-3}{2}{-2}
\mcq{382}{The equation y = 3x represents which type of graph:}{A line with slope 3}{A parabola}{A cubic curve}{Exponential decay}
\mcq{383}{The gradient of the tangent to the curve y = 2x$^{2}$ - 3x + 1 at x = 1:}{1}{2}{3}{4}
\mcq{384}{The equation y = - 2x$^{2}$+ 4x - 3 represents which type of graph:}{Parabola opening downwards}{Parabola opening upwards}{Line}{Exponential growth}
\mcq{385}{In a reciprocal function $y=\dfrac{1}{x}$ as x approaches zero, what happens to y?}{It approaches infinity}{It approaches zero}{It becomes undefined}{It remains constant}
\mcq{386}{The general form of a quadratic function is:}{y = mx + c}{y = ax$^{2}$+ bx + c}{y = ax$^{3}$+ bx$^{2}$+ cx + d}{y = k x e$^{2}$x}
\mcq{387}{In the function y = x$^{2}$ + 4x + 4 what is the value of y when x = - 2:}{0}{4}{-4}{8}
\mcq{388}{In the equation y = 5x - 2 , what is the y-intercept:}{-5}{-2}{2}{5}
\mcq{389}{What is the break-even point when C(x) = 1200 + 20x and R(x) = 50x:}{40}{50}{60}{75}
\mcq{390}{In the function y = x$^{3}$+ 3x$^{2}$ what is the slope of the tangent at x = 0:}{0}{3}{6}{-3}
\mcq{391}{The function y = 3x$^{2}$+ 2x - 1 has its minimum value at:}{x = 0}{x = -1}{x = 1}{x = 2}
\mcq{392}{In the function y = x$^{2}$+ 4x what is the axis of symmetry:}{x = 0}{x = -2}{x = 2}{x = -4}
\mcq{393}{Which of the following is the domain of the reciprocal function (eq):}{All real numbers except zero}{All real numbers}{Positive real numbers}{Negative real numbers}
\mcq{394}{The graph of y = 2x$^{3}$- x$^{2}$- 5 is:}{Linear}{Quadratic}{Cubic}{Exponential}
\mcq{395}{The graph of y = 2x$^{3}$- x$^{2}$- 5 is:}{Linear}{Quadratic}{Cubic}{Exponential}
\mcq{396}{For the function y = 3x - 4 the graph is a:}{Line}{Parabola}{Hyperbola}{Exponential function}
\mcq{397}{The graph of y = x$^{3}$ passes through the point:}{(0,0)}{(1,0)}{(0,1)}{(0,-1)}
\mcq{398}{What is the y-intercept of the equation y = 4x - 6:}{-6}{6}{4}{-4}
\mcq{399}{The cost function is C(x) = 1500 + 10x The revenue function is R(x) = 50x What is the break-even point:}{30 units}{50 units}{100 units}{150 units}
\mcq{400}{The graph of y = x$^{2}$- 2x + 1 represents:}{A line}{A parabola}{A cubic curve}{An exponential curve}
\mcq{401}{In the function y = 2x$^{2}$+ 3x + 5 , the value of y when x = - 1 is:}{4}{2}{1}{3}
\mcq{402}{The graph of y = x$^{2}$- 4 intersects the x-axis at:}{(2, 0) and (-2, 0)}{(-4, 0)(4,0) and (-4, 0)}{(1,0) and (-1,0)}{(3, 0) and (-3, 0)}
\mcq{403}{The function y = x$^{4}$ is a:}{Linear function}{Quadratic function}{Cubic function}{Quartic function}
\mcq{404}{The graph of y = x$^{2}$ + x - 2 represents:}{A straight line}{A hyperbola}{A parabola}{An exponential function}
\mcq{405}{For the function y = x$^{2}$ + 2x + 1, what is the vertex of the parabola:}{(0,0)}{(-1,0)}{(-1,1)}{(1,1)}
\mcq{406}{The graph of y = x$^{2}$ + 5 intersects the x-axis at:}{(0,5)}{(5,0)}{(-5,0)}{(1,0)}
\mcq{407}{The function y = 3x$^{2}$ represents:}{A line}{A parabola opening upwards}{A parabola opening downwards}{A cubic function}
\mcq{408}{The function y = 3x$^{2}$ represents:}{A line}{A parabola opening upwards}{A parabola opening downwards}{A cubic function}
\mcq{409}{The function y=2$^{x}$ represents what kind of growth:}{Linear growth}{Exponential growth}{Exponential decay}{Quadratic growth}
\mcq{410}{What is the gradient of the tangent to the curve y=x$^{2}$ at x=1:}{4}{2}{0}{8}
\mcq{411}{Which function represents exponential decay:}{y=2$^{x}$}{y=e$^{x}$}{y=2$^{-x}$}{y=x$^{2}$}
\mcq{412}{In the function y=x$^{2}$-3x-2, what is the value of y when x=0:}{-2}{0}{2}{-3}
\mcq{413}{The graph of the function y=x$^{2}$-2 represents a:}{Parabola opening upwards}{Parabola opening downwards}{Hyperbola}{Circle}
\mcq{414}{What is the break-even point when the cost function is C(x)=90,000+600x and the reverse function is R(x)=1,800x:}{50 footballs}{75 footballs}{100 footballs}{150 footballs}
\mcq{415}{Which of the following is an exponential growth function:}{y = x$^{2}$}{y = 3x}{y = 2$^{x}$}{y = x$^{3}$}
\mcq{416}{Which of the following represents an exponential decay function:}{y = 2$^{x}$}{y = e$^{-x}$}{y = x$^{2}$}{y = 3x}
\mcq{417}{The graph of y = e$^{x}$ represents:}{Linear growth}{Exponential growth}{Parabolic growth}{Cubic growth}
\mcq{418}{x = 5 represents:}{x-axis}{y-axis}{line $\parallel$ to x$-$axis}{line $\parallel$ to y$-$axis}
\mcq{419}{Slope of the line y = 5x + 3 is:}{3}{-3}{5}{-5}
\mcq{420}{The y-intercepts of y= -2x - 1 is:}{-2}{2}{-1}{1}
\mcq{421}{The graph of $y=x^{3}$, cut the x-axis at:}{x=0}{x=1}{x=-1}{x=2}
\mcq{422}{The graph of $3^{x}$, cuts the x-axis at:}{Growth}{Decay}{Both}{A line}
\mcq{423}{The graph of $y=-x^{2}+5$ :}{Upward}{Downward}{Left side}{Right side}
\mcq{424}{The graph of $y=x^{2}-9$ :}{Upward}{Downward}{Left side}{Right side}
\mcq{425}{y=5$^{x}$ is ...... function:}{Linear}{Quadratic}{Cubic}{Exponential}
\mcq{426}{Reciprocal function is:}{$y=7^{x}$}{$y=\dfrac{2}{x}$}{$y=2x^{2}$}{$y=5x^{3}$}
\mcq{427}{$y=-3x^{3}+7$ is ...... function:}{Exponential}{Cubic}{Linear}{Reciprocal}
\mcq{428}{The point of intersection of the perpendicular bisectors of the sides of a triangle is called:}{Incenter}{Centroid}{Circumcenter}{Orthocenter}
\mcq{429}{If the measure of one angle in a triangle is greater than 90$^{\circ}$, which type of triangle is formed:}{Acute triangle}{Right-angled triangle}{Obtuse triangle}{Isosceles triangle}
\mcq{430}{The point where all the medians of a triangle meet is known as:}{Circumcenter}{Incenter}{Centroid}{Orthocenter}
\mcq{431}{Which locus defines a point that is a fixed distance from a given point:}{Perpendicular bisector}{Circle}{Angle bisector}{Line segment}

\mcq{432}{What is the angle bisector of a triangle:}{A line passing through a vertex}{A line that divides an angle into two equal parts}{A line joining midpoints of sides}{A line perpendicular to a side}
\mcq{433}{The orthocenter of a triangle is where:}{The altitudes of the triangle meet}{The sides are perpendicular}{The circumcenter meets the centroid}{None of these}
\mcq{434}{The intersection of the three angle bisectors of a triangle is called:}{Orthocenter}{Centroid}{Circumcenter}{Incenter}
\mcq{435}{An equilateral triangle is .........................angled triangle:}{Obtuse}{Right}{Acute}{None of these}
\mcq{436}{The center of a circle that passes through all vertices of a triangle is located at the:}{Incenter}{Circumcenter}{Centroid}{Orthocenter}
\mcq{437}{How many types of a triangle w.r.t angles:}{2}{3}{4}{5}
\mcq{438}{A triangle with all sides and angles equal is called a:}{Right triangle}{Isosceles triangle}{Equilateral triangle}{Scalene triangle}
\mcq{439}{The centroid of a triangle is the point of concurrency of the:}{Angle bisectors}{Altitudes}{Perpendicular bisectors}{Medians}
\mcq{440}{When constructing a right-angled triangle, the side opposite the right angle is called the:}{Hypotenuse}{Median}{Altitude}{Base}
\mcq{441}{Which of the following is true for an equilateral triangle:}{It has only one angle of 90$^{\circ}$}{All angles are 60$^{\circ}$}{It has two equal sides}{The sides are all different}
\mcq{442}{What is the sum of the interior angles of any triangle:}{180$^{\circ}$}{360$^{\circ}$}{90$^{\circ}$}{120$^{\circ}$}
\mcq{443}{A triangle with two equal sides is called a:}{Equilateral triangle}{Scalene triangle}{Isosceles triangle}{Right triangle}
\mcq{444}{How many types of a triangle w.r.t sides:}{2}{3}{4}{5}
\mcq{445}{The centroid of a triangle divides each median into two parts in a ratio of:}{1:2}{2:1}{1:3}{3:1}
\mcq{446}{Which of the following is not a type of triangle based on sides:}{Scalene}{Isosceles}{Equilateral}{Acute}
\mcq{447}{Perpendicular from a vertex to the opposite side in a triangle is called:}{Altitude}{Median}{Bisector}{Angle}
\mcq{448}{A triangle with one angle equal to 90$^{\circ}$ is called a:}{Scalene triangle}{Equilateral triangle}{Right-angled triangle}{Isosceles triangle}
\mcq{449}{A triangle is a closed figure having:}{Two sides}{Three sides}{Four sides}{Five sides}
\mcq{450}{When constructing a triangle, which condition requires the measure of one angle and two sides:}{SAS criterion}{ASA criterion}{SSS criterion}{AAS criterion}
\mcq{451}{What is the angle formed by two intersecting rays from a vertex in a triangle called:}{Complementary angle}{Opposite angle}{Vertical angle}{Exterior angle}
\mcq{452}{A triangle can be constructed if the sum of the measure of the any two sides is ...... the measure of the third side:}{Less than}{Greater than}{Equal to}{Greater than and equal to}
\mcq{453}{An equilateral triangle ......:}{Can be isosceles}{Can be right angled}{Can be obtuse angled}{Has each angle equal to 50$^{\circ}$}
\mcq{454}{If the sum of the measure of two angles is less than 90$^{\circ}$. then the triangle is ......:}{Equilateral}{Acute angled}{Obtuse angled}{Right angled}
\mcq{455}{The line segment joining the midpoint of a side to its opposite vertex in a triangle is called ......:}{Median}{Perpendicular bisector}{Angle bisector}{Circle}
\mcq{456}{The angle bisectors of a triangle intersect at ......:}{One point}{Two points}{Three points}{Four points}
\mcq{457}{Locus of all points equidistant from a fixed point is ......:}{Circle}{Perpendicular bisector}{Angle bisector}{Parallel lines}
\mcq{458}{Locus of points equidistant from two fixed points is ......:}{Circle}{Perpendicular bisector}{Angle bisector}{Parallel lines}
\mcq{459}{Locus of points equidistant from a fixed line is/are ......:}{Circle}{Perpendicular bisector}{Angle bisector}{Parallel lines}
\mcq{460}{Locus of points equidistant from two intersecting lines is ......:}{Circle}{Perpendicular bisector}{Angle bisector}{Parallel lines}
\mcq{461}{The set of all points which is farther than 2 km from a fixed point B is a region outside a circle of radius ...... and centre at B:}{1 km}{1.9 km}{2 km}{2.1 km}
\mcq{462}{What is the process of collecting, organizing, summarizing, analyzing, and interpreting numerical data called:}{Data Collection}{Information Handling}{Data Analysis}{Data Interpretation}
\mcq{463}{Which of the following is an example of discrete data:}{Height of students}{Weight of students}{Number of books sold}{Temperature of a city}
\mcq{464}{What is the term for data that can take every possible value in a given interval:}{Discrete data}{Continuous data}{Grouped data}{Ungrouped data}
\mcq{465}{Which of the following is an example of ungrouped data:}{Marks of students arranged in classes}{Number of toys sold in a month}{Height of students in a class}{Temperature recorded daily}
\mcq{466}{What is the term for data arranged in systematic order (groups or classes):}{Ungrouped data}{Grouped data}{Discrete data}{Continuous data}
\mcq{467}{How used the word data handling first:}{Netwon}{Henry}{Sir Rondddfisher}{Arther calay}
\mcq{468}{What is the formula to calculate the size of class limits:}{Size = Range / Number of classes}{Size = Range x Number of classes}{Size = Range + Number of classes}{Size = Range - Number of classes}
\mcq{469}{What is the range of the data if the greatest value is 136 and the smallest value is 30:}{106}{166}{96}{136}
\mcq{470}{What is the term for the difference between the upper and lower class limits:}{Class size}{Class boundary}{Class interval}{Class frequency}
\mcq{471}{In class 10-14 lower class limits is:}{9.5}{10}{12}{146}
\mcq{472}{What is a histogram:}{A graph of adjacent rectangles}{A graph of connected lines}{A graph of pie charts}{A graph of bar charts}
\mcq{473}{What is the term for the midpoint of a class interval:}{Class boundary}{Class mark}{Class limit}{Class frequency}
\mcq{474}{What is the formula to calculate the midpoint of a class interval?}{$Midpoint=\dfrac{Lower\ time\ limit+Upper\ time\ limit}{2}$}{$Midpoint=\dfrac{Lower-Upper\ time}{2}$}{$Midpoint=\dfrac{Lower\ limit\times Upper\ limit}{2}$}{None of these}
\mcq{475}{In class 10-14 frequency is 10 then written in totally number:}{\tb\tb}{\tb\tb\tb\tb\tb}{\tb\tb\tb\tb}{None}
\mcq{476}{Class Boundaries of 5-9 is:}{5.5-10.5}{4.5-8.5}{4.5-9.5}{4-10}
\mcq{477}{Mode of first 20 whole numbers:}{0}{20}{19}{No Mode}
\mcq{478}{What is the median of the data set 5, 8, 12, 13, 16, 18, 20:}{12}{13}{16}{18}
\mcq{479}{What is the mode of the data set 10, 15, 15, 20, 20, 20, 25:}{15}{20}{25}{No mode}
\mcq{480}{What is the term for the most frequently occurring value in a data set:}{Mean}{Median}{Mode}{Range}
\mcq{481}{What is the formula for the weighted mean?}{Weighted=$\dfrac{\sum(wx)}{\sum w}$}{Weighted mean =(wx) $\times$ Ew}{Weighted mean =(wx) $-$ S}{Weighted mean = (wx) + Ew}
\mcq{482}{What is the weighted mean if the weights are 4, 3, 5, 6 and the values are 74, 78, 74, 90:}{80}{85}{90}{95}
\mcq{483}{What is the term for the average where different observations have different weights:}{Arithmetic mean}{Weighted mean}{Median}{Mode}
\mcq{484}{What is the average number of packets of medicine distributed by a company if the total number of packets is 114 and the number of areas is 7:}{16}{17}{18}{19}
\mcq{485}{What is the average budget for the next year if the yearly budgets are 5, 7, 8, 6, 7 million:}{6.6}{7.0}{7.5}{8.0}
\mcq{486}{Arithmetic mean denoted by:}{$\widetilde{X}$}{$\overline{X}$}{$\widehat{X}$}{SD}
\mcq{487}{What is the average marks of a student if the marks are 84, 91, 72, 68, 87, 78:}{80}{82}{84}{86}
\mcq{488}{Symbol used for sum:}{$\equiv$}{$\sum$}{\%}{$\sim$}
\mcq{489}{What is the median of the data set 55, 53, 54, 58, 60, 61, 62, 56, 57, 51, 63:}{56}{57}{58}{59}
\mcq{490}{Formula used for arithmetic mean for ungroup:}{$\dfrac{n+1}{2}$ th term}{$\dfrac{\sum x}{n}$}{Both a and b}{None of these}
\mcq{491}{What is the median of the data set 110, 125, 122, 130, 124, 127, 120:}{124}{125}{126}{127}
\mcq{492}{What is the mode of the data set 88, 70, 72, 125, 95, 81, 90, 95, 90:}{90}{95}{90 and 95}{No mode}
\mcq{493}{Symbol $\widetilde{X}$ used for:}{Mode}{Median}{Range}{Arithmetic mean}
\mcq{494}{What is the arithmetic mean of the data set 4, 6, 10, 12, 15, 20, 25, 28, 30:}{16.67}{18.00}{20.00}{22.00}
\mcq{495}{What is the arithmetic mean of the data set 12, 18, 19, 0, -19, -18, -12:}{0}{1}{2}{3}

\mcq{496}{$l+\dfrac{f_{m}-f_{1}}{(f_{m}-f_{1})+(f_{m}-f_{2})}\times h$ used for:}{Mode}{Median}{Range}{Arithmetic mean}
\mcq{497}{Which data takes only some specific values:}{Continuous data}{Discrete data}{Grouped data}{Ungrouped data}
\mcq{498}{The number of times a value occurs in a data is called:}{Frequency}{Relative frequency}{Class limit}{Class boundaries}
\mcq{499}{Midpoint is also known as:}{Mean}{Median}{Class limit}{Class mark}
\mcq{500}{Frequency polygon is also drawn / constructed by using:}{Histogram}{Bar graph}{Class boundaries}{Class limit}
\mcq{501}{The difference between the greatest value and the smallest value is called:}{Class limit}{Midpoint}{Relative frequency}{Range}
\mcq{502}{Measure of central tendency is used to find out the ...... of a data set:}{Class boundaries}{Cumulative frequency}{Middle or centre value}{Frequency}
\mcq{503}{If the mean of 5,7,8,9 and x is 7,5, what will be the value of x:}{10}{8}{8.5}{5.8}
\mcq{504}{Find the mode of the given data 2,5,8,9,0,1,3,7 and 10:}{5}{7}{0}{No mode}
\mcq{505}{In a data the values (observation) which appears or occurs most often is called:}{Mean}{Mode}{Median}{Weighted mean}
\mcq{506}{Find the median of the given data 110,125,122,130,124 and 120:}{124}{120}{125}{127}
\mcq{507}{What is the formula to calculate probability:}{$P(A)=\dfrac{n(A)}{n(S)}$}{$P(A)=\dfrac{n(S)}{n(A)}$}{P(A) = n(A) x n(S)}{P(A) = n(S) $-$ n(A)}
\mcq{508}{Which of the following is not an example of an experiment:}{Tossing a coin}{Rolling a dice}{Guessing a number}{Drawing a card from a deck}
\mcq{509}{What is the probability of getting a head when tossing a fair coin?3.}{$\dfrac{1}{2}$}{$\dfrac{1}{3}$}{$\dfrac{1}{4}$}{$\dfrac{1}{6}$}
\mcq{510}{If a fair dice is rolled, what is the probability of getting a number divisible by 3?}{$\dfrac{1}{6}$}{$\dfrac{1}{4}$}{$\dfrac{1}{3}$}{$\dfrac{1}{2}$}
\mcq{511}{What is the sample space when a single coin is tossed:}{\{A, B\}}{\{1, 2, 3\}}{\{1, 2, 3, 4\}}{\{H, T\}}
\mcq{512}{If a dice is rolled twice, what is the total number of possible outcomes:}{6}{36}{12}{24}
\mcq{513}{If event A has 3 favorable outcomes, and the sample space contains 6 possible outcomes, what is P(A) :}{$\dfrac{1}{6}$}{$\dfrac{1}{4}$}{$\dfrac{1}{3}$}{$\dfrac{1}{2}$}
\mcq{514}{The complement of event A is:}{P(S)}{P(A)}{P(A$'$)}{P(A and B)}
\mcq{515}{What is the probability of not getting a 6 on a fair dice?}{$\dfrac{1}{6}$}{$\dfrac{5}{6}$}{$\dfrac{1}{3}$}{$\dfrac{2}{3}$}
\mcq{516}{If a fair dice is rolled, what is the probability of getting an even number ?}{$\dfrac{1}{2}$}{$\dfrac{1}{3}$}{$\dfrac{1}{4}$}{$\dfrac{1}{6}$}
\mcq{517}{When a fair dice is rolled twice, what is the probability of getting a sum of 8:}{$\dfrac{1}{9}$}{$\dfrac{5}{36}$}{$\dfrac{1}{36}$}{$\dfrac{1}{12}$}
\mcq{518}{What is the probability of getting a double six when rolling two dice:}{$\dfrac{1}{36}$}{$\dfrac{1}{18}$}{$\dfrac{1}{12}$}{$\dfrac{1}{6}$}
\mcq{519}{The probability of an event is always:}{Less than 0}{Greater than 1}{Equal to 1}{Between 0 and 1}
\mcq{520}{The probability of selecting a blue ball from a bag of 5 blue and 8 green balls is:}{$\dfrac{5}{13}$}{$\dfrac{8}{13}$}{$\dfrac{1}{3}$}{$\dfrac{5}{8}$}
\mcq{521}{What is the probability of getting neither a spade nor a heart when drawing a card from a deck of 52?}{$\dfrac{26}{52}$}{$\dfrac{26}{39}$}{$\dfrac{1}{2}$}{$\dfrac{1}{4}$}
\mcq{522}{If the probability of an event is 0.3, the probability of its complement is:}{0.4}{0.5}{0.3}{0.7}
\mcq{523}{What is the probability of drawing a red heart from a deck of 52 playing cards?}{$\dfrac{1}{13}$}{$\dfrac{1}{4}$}{$\dfrac{1}{2}$}{$\dfrac{1}{26}$}
\mcq{524}{What is the probability of not getting a red ball from a bag containing red and 5 blue balls?}{$\dfrac{5}{13}$}{$\dfrac{8}{13}$}{$\dfrac{5}{8}$}{$\dfrac{8}{5}$}
\mcq{525}{If 3 dice are rolled, what is the probability of getting at least one 6:}{$\dfrac{1}{36}$}{$\dfrac{3}{36}$}{$\dfrac{1}{6}$}{$\dfrac{5}{36}$}
\mcq{526}{The probability of getting even sum when rolling two dice is:}{$\dfrac{5}{36}$}{$\dfrac{18}{36}$}{$\dfrac{9}{36}$}{$\dfrac{12}{36}$}
\mcq{527}{The probability of drawing a red card from a deck of 52 cards is:}{$\dfrac{1}{2}$}{$\dfrac{13}{52}$}{$\dfrac{26}{52}$}{$\dfrac{1}{26}$}
\mcq{528}{What is the probability of drawing a white ball from a bag containing 5 white and 7 black balls?}{$\dfrac{1}{3}$}{$\dfrac{7}{12}$}{$\dfrac{1}{2}$}{$\dfrac{5}{12}$}
\mcq{529}{The relative frequency is:}{Total number of favorable outcomes}{Total number of possible outcomes}{How often an event occurs}{The number of trials}
\mcq{530}{The relative frequency of selecting a red ball from 80 students is:}{0.29}{0.19}{0.31}{0.12}
\mcq{531}{If six dice are rolled 50 times, the expected frequency of a number appearing is given by:}{NP(B)}{NP(A)}{NP(X)}{None of these}
\mcq{532}{The relative frequency of green in a survey of 80 students is:}{$\dfrac{15}{80}$}{0.19}{0.31}{$\dfrac{10}{80}$}
\mcq{533}{If a probability is 1/5, the expected frequency in 100 trials is:}{20}{50}{10}{25}
\mcq{534}{Each element of the sample space is called:}{Event}{Experiment}{Sample point}{Outcomes}
\mcq{535}{An outcome which represents how many times we expect something to happen is called:}{Outcomes}{Favourable outcome}{Sample space}{Sample point}
\mcq{536}{Which one tells us how often a specific event occurs relative to the total number of events:}{Expected frequency}{Sum of relative frequency}{Relative frequency}{Frequency}
\mcq{537}{Estimated probability of an event occurring is also known as:}{Relative frequency}{Expected frequency}{Class boundaries}{Sum of expected frequency}
\mcq{538}{The sum of all expected frequencies is equal to the fixed number of:}{Trials}{Relative frequencies}{Outcomes}{Events}
\mcq{539}{Find out the total number of possible sample space when 4 dice are rolled:}{6$^{2}$}{6$^{3}$}{6$^{4}$}{6$^{6}$}
\mcq{540}{The chance of occurrence of a particular event is called:}{Sample space}{Estimated probability}{Probability}{Expected frequency}
\mcq{541}{While rolling a pair of dice, what will be the probability of double 2?}{$\dfrac{1}{6}$}{$\dfrac{1}{3}$}{$\dfrac{5}{6}$}{$\dfrac{1}{36}$}
\mcq{542}{A card is chosen from a pack of 52 playing cards, find the probability of getting no jack and king?}{$\dfrac{2}{13}$}{$\dfrac{11}{13}$}{$\dfrac{2}{52}$}{$\dfrac{11}{52}$}
\mcq{543}{An event which will probably occur. It has greater chance to occur is called:}{Equally likely event}{Likely event}{Unlikely event}{Certain event}


\end{document}
